\subsection{From a Cap to an SLO}
\label{sec:design-slo}

A cumulative cap alone produces the wrong behaviour at both ends: a tenant that
spends a subject's budget in the first hour gets full service for an hour and
nothing thereafter, while one that spreads the same spend over thirty days gets
uniform service. Neither is what the subject agreed to. \sys\ therefore evaluates a population-level SLO:
over a window of length $W$, the fraction of monitored subjects whose spend
exceeds $\theta_{\mathrm{soft}}=\alpha E$ must not exceed $p$, and no subject may
exceed $E$. The default is $\alpha=0.8$, $p=0.01$, $W=600$\,s and $\ell=300$\,s
epochs, so the ring holds three epochs.

The monitor reads the windowed estimates of the hot table rather than of the whole
population, on the argument that a subject outside the hot set has, by construction,
accumulated less than the eviction threshold and cannot be the subject whose spend
violates the soft threshold. This keeps the monitor's cost independent of $N$ and
concentrates attention where violations are possible, but the reported rate is over
the \emph{violation-capable} subpopulation rather than over all subjects. The
tenant-level rate is refreshed every 64 evaluations, bounding staleness at the cost
of a lag between a change in behaviour and the breaker's response.

The breaker is a four-level machine (normal, throttle, reject, freeze) that
escalates immediately and de-escalates one level at a time. De-escalation
requires pressure below the current threshold by a margin $h$ for $d$ consecutive
evaluations: the margin prevents oscillation around a threshold, the dwell prevents
one quiet evaluation from releasing a subject about to be charged again. In throttle
mode the charge is scaled by $\beta=0.5$, which inflates the noise by
$1/\sqrt{\beta}$; the tenant receives a lower-quality answer at a lower price rather
than a refusal. This turns a hard wall into a degradation curve, and it is sound
because the mechanism executed is the one charged for.

\subsection{Two-Phase Settlement}
\label{sec:design-settle}

Reservation is all-or-nothing over the query's subject set. Under a guard covering
exactly the stripes of those subjects, the coordinator reads
$\mathrm{remaining}(s)=E_s-(\text{committed}(s)+\text{in-flight}(s))$ for every
$s$ in the support and rejects if any subject lacks room; otherwise it records a
reservation with a time-to-live. Check and record share the guard, so no other
transaction observes the budget between them, and a rejection is not a partial
charge: either the query is reserved against every subject it touches, or it
touches no budget at all.

Commit is idempotent in the transaction identifier. A repeated commit for an
already-committed transaction increments a duplicate counter and returns the
existing record rather than charging again, which makes client retries safe.
Release returns the reserved budget and is a no-op for a transaction already
committed or released. Reservations neither committed nor released within their
TTL are expired by a periodic tick, bounding the budget a crashed client can hold
hostage to one TTL; the cost is a window in which budget is unavailable although
no answer was delivered, so the TTL trades availability under crash against
availability under load.

Concurrency is handled by striping. Subjects map to $L=64$ stripes by a hash of
the identifier, and a transaction acquires its subjects' stripes in sorted index
order, which rules out deadlock and lets transactions over disjoint subjects
proceed in parallel. A transaction spanning many subjects serialises against everything sharing its
stripes, which is correct for a multi-subject query and is why striped and global
locking behave alike when all traffic hits one subject. The difference appears
when subjects are disjoint, and Section~\ref{sec:eval-concurrency} measures
both.
